\section{A parameter-free modus-ponens command using a retained gate residual}
\label{sec:residual-parameter}

This section proves completeness of an exact revised command language.  Its
purpose is to remove the numeric input field of the MP command.  It is a
semantic theorem about the command interpreter below, not by itself a claim
about the number of states of a compiled binary machine.  The literal source
splice, compiler cost, and binary runtime obligations must be reported
separately.

The useful observation is that a number discarded by AX's final freshness
comparison can serve as the consequent parameter of a later MP.  It need not
be the original AX argument.  A surjective residual suffices, because a dummy
AX can set the parameter while the actual proof premises are protected on the
theorem stack.

\subsection{Exact command interpreter}

Write a work state as $(t,S,d,h)$: current theorem code, theorem stack,
numeric parameter, and cached theorem code.  Stacks use the same Cantor
pairing as the accepted machine, with $[x]S=P(x,S)$.  A displayed finite list
is pushed with its leftmost entry at the top.  The empty stream convention is
$P^{-1}(0)=(0,0)$.  Write $I(a,b)=4P(a,b)$, $Q_i(a)=4P(i,a)+2$, and
$F=1$.

Let $\rho_1,\ldots,\rho_8$ denote the unchanged native numerical rows;
the seventh entry is the repository's \nolinkurl{growG} row.  Only $\rho_8$
uses the fourth input.  Let
\[
 g(a,b,c)=
 \begin{cases}
  0,&c=a\text{ or }c=b,\\
  \operatorname{bundle}(a,b,c),&\text{otherwise},
 \end{cases}
 \qquad
 k(c,a)=\max\{a-c-1,0\}.
\]
Every $\rho_j$, every $Q_i(a)$, and every value of $g$ differs from $1$.
This is an exact constructor-tag fact, independent of derivability: the
rows and nonzero bundle are implication codes, and $Q$ has tag two.

The new commands have the following complete effects:
\begin{align*}
 \mathsf{DROP}(t,S,d,h)&=(h,[t]S,d,h),\\
 \mathsf{NIP}(t,[x]S,d,h)&=
  \begin{cases}\mathsf{halt},&h=1,\\(t,S,d,x),&h\ne1,\end{cases}\\
 \mathsf{MP}(t,S,d,h)&=
  \begin{cases}(d,[0]S,d,0),&I(t,d)=h,\\(0,S,d,0),&I(t,d)\ne h.\end{cases}
\end{align*}
In particular, none of these three commands reads a numeric field.
For $\mathsf{AX}(c,y,x)$, put $a=x$, $b=h+y$, and $u=d$, using the
\emph{old} values of $d,h$.  Its effect is
\[
 (t,S,d,h)\longmapsto
 \bigl(g(a,b,c),
 [\rho_1(a,b),\ldots,\rho_8(a,b,c,u),Q_u(h),t]S,
 k(c,a),0\bigr).
\]
The suppressed arguments of $\rho_2,\ldots,\rho_7$ are exactly those
of the accepted emitter.  There are nine new theorem entries and then the
old current code, so AX pushes exactly ten entries.

This interpreter emits unrestricted $\rho_8$: the old parameter $u$ and the
new absolute input $a$ are independent.  That is sound for the original
native calculus.  It does not rely on the dominated-row certificate.

\paragraph{Opposite-residual variant.}
The same command-calibration theorem holds if AX instead returns
\[
 k_X(c,a)=\max\{c-a-1,0\}.
\]
In that variant replace the setter below by
\[
 \mathsf{SET}_X(v)=
 \mathsf{AX}(v+1,v+1,0);N^{10};\mathsf{DROP};\mathsf{NIP}.
\]
For incoming cache zero this has $a=0$ and $b=c=v+1$, hence gate zero and
$k_X(c,a)=v$.  The proof that it preserves the current theorem and old
stack is identical.  RESET and the generalization AX use $a=c=0$, so both
residuals are zero there; the leaf and bundle contracts leave the output
parameter unconstrained.  Every subsequent compiler argument therefore
applies with this single setter replacement.  This variant corresponds
to retaining the scratch-counter remainder, whose transfer to $D$ can
reuse the source's existing scratch-to-parameter transfer macro.

\subsection{Why this residual is available at the source level}

In the accepted source the gate first consumes the old cache while comparing
$b$ with $c$.  It then transfers $a$ from register $D$ to the now-zero cache
register $H$ and performs the second comparison.  In the optimized K217
comparison it repeatedly decrements $H$ first and then its scratch copy of
$c$.  If $a>c$, the first exhausted scratch check leaves $a-c-1$ in $H$;
if $a\le c$, $H$ is zero before the unequal cleanup or equality branch.
The unequal cleanup ordinarily clears both remainders and skips the branch
that clears the current bundle.  Replace \emph{only its clearing loop for
$H$ in this second comparison} by the transfer $H\to D$.  Then the
comparison still clears the bundle exactly when $c=a$, its scratch and cache
registers still finish at zero, and $D$ finishes at $k(c,a)$.  The first
comparison and its clearing behavior remain unchanged.

Here is the exhaustive case analysis for the optimized comparison.  If
$a=c$, the paired decrements exhaust both counters together, the equality
branch clears the current bundle, and $D$ stays zero.  If $a>c$, the $c$
paired iterations are followed by one additional decrement of $H$; its
remaining $a-c-1$ units are transferred to $D$, and the current bundle is
preserved.  If $c>a$, the $a$ paired iterations are followed by one
additional decrement of the scratch counter; $H$ is already zero, so the
transfer is empty, the scratch remainder is cleared, and the current bundle
is preserved.  These cases prove the residual contract for arbitrary
natural inputs, beyond the finite numerical tests.

The other proposed local change makes the AX input $a$ absolute.  After GEN
has used the old parameter as its quantifier label, the old parameter is
loaded for the subformula $C\mathbin{\to}U$ of $\rho_8$.  Consume it on
this final use instead of copying and restoring it.  The subsequent additive
read into $D$ then starts at zero, and therefore reads $a=x$.  All formula
values emitted before and after this load have the values in the displayed
AX contract.  These statements specify the source obligations; they must be
checked against the literal placement of the relevant macros, rather than
inferred just from register names.

\subsection{The exact helper blocks}

Let $N^j$ mean $j$ repetitions of $\mathsf{NIP}$.  Command sequences are
written from left to right in execution order.  Early halting is acceptable
for a proof block: the purpose of the full interpreter is to halt precisely
when a refutation exists.

\paragraph{Reset.}
Define
\[
 \mathsf{RESET}
  =\bigl(\mathsf{AX}(0,0,0);N^{10}\bigr)^2.
\]
For \emph{every} numerical work state, including arbitrary initial cache
and current values,
\[
 \mathsf{RESET}(t,S,d,h)=(0,S,0,0).
\]
There is no early halt in this block.  In its first half AX has $a=c=0$,
so its gate and residual are both zero.  The ten NIPs successively pop nine
entries known not to equal $1$, followed by the old $t$.  The latter is not
tested by that last NIP: the result is $(0,S,0,t)$.  The second AX resets
the cache before another NIP can inspect it.  The same calculation now pops
an old current value of zero and gives $(0,S,0,0)$.  This distinction between
the cache checked on entry and the entry popped on exit is essential.

\paragraph{Parameter setter preserving the current theorem.}
For any target natural number $v$, define
\[
 \mathsf{SET}(v)=
 \mathsf{AX}(0,0,v+1);N^{10};\mathsf{DROP};\mathsf{NIP}.
\]
Starting with cache zero,
\[
 \mathsf{SET}(v)(t,S,d,0)
  =\mathsf{halt}\quad\text{or}\quad(t,S,v,0).
\]
More precisely, this block halts exactly when $t=1$.  The first gate is zero
because $b=c=0$, and the residual is
$k(0,v+1)=v$.  After ten NIPs the state is $(0,S,v,t)$.  DROP changes it to
$(t,[0]S,v,t)$; the final NIP either halts on $t=1$, or pops the pushed zero
and restores $(t,S,v,0)$.  Both the current theorem and every old stack entry
are preserved.  In particular, an arbitrary implication can safely be held
current while this setter runs, because implication codes differ from $1$.

\paragraph{Selecting an emitted row.}
Suppose AX was entered with current value zero and returns
\[
 (g,[r_1,\ldots,r_9,0]S,d',0),\qquad
 g\ne1,\quad r_1,\ldots,r_9\ne1.
\]
For $1\le j\le9$, the block
\[
 N^j;\mathsf{DROP};N^{11-j}
\]
returns exactly $(r_j,S,d',0)$.  The first $j$ NIPs put $r_j$ in the cache.
DROP makes it current and pushes $g$.  The remaining stack prefix consists
of $g$, the $9-j$ following rows, and the old zero current value, a total
of $11-j$ entries.  All cache values tested during these pops differ from
$1$.  This proves both the offset and the exact restoration of the old
stack.

\subsection{Compilation of every native derivation}

Call a finite block $p$ a compiler block for $X$ when, for every
$(t,S,d,h)$, execution either halts or returns $(X,S,d',0)$ for some
$d'$.  There is no precondition on any of the four initial values.
The old stack is restored exactly; the numeric output parameter need not
be zero.  The following induction constructs such a block for every native
derivable code.

\paragraph{Truth.}
Use $\mathsf{RESET}$, whose exact contract was proved above.

\paragraph{Native row $j$, $1\le j\le8$.}
To obtain the row at parameters $(a,b,c,u)$, use
\[
 \mathsf{RESET};\mathsf{SET}(u);\mathsf{AX}(c,b,a);
 N^j;\mathsf{DROP};N^{11-j}.
\]
After RESET the current value and cache are zero.  SET therefore finishes
without halting and supplies the desired old parameter $u$.  AX reads
$a,b,c$ independently and emits the desired row.  The selection contract
finishes with that row current, the original stack restored, and cache zero.
Unused parameters of the earlier rows can be chosen arbitrarily.

\paragraph{Native bundle.}
When $c\ne a,b$, use
\[
 \mathsf{RESET};\mathsf{AX}(c,b,a);N^{10}.
\]
The gate equals the requested bundle, and the ten safe pops remove the
nine emitted rows and the old current zero.  The bundle remains current,
the old stack is restored, and the cache is zero.

\paragraph{Generalization.}
If $p$ is a compiler block for $H$, a compiler block for $Q_i(H)$ is
\[
 p;\mathsf{SET}(i);\mathsf{DROP};\mathsf{NIP};
 \mathsf{AX}(0,0,0);N^9;\mathsf{DROP};N^2.
\]
If $p$ or SET halts, the required alternative already holds.  Otherwise the
state after SET is $(H,S,i,0)$.  DROP and NIP give $(0,S,i,H)$.
The next AX emits $Q_i(H)$ in its ninth stack position and has gate zero,
because $a=c=0$.  Its old current entry is zero.  The row-selection contract
with $j=9$ gives $(Q_i(H),S,0,0)$.

\paragraph{Modus ponens.}
If $p_A$ and $p_{I(A,B)}$ compile the two premises, use
\[
 p_A;\mathsf{DROP};p_{I(A,B)};\mathsf{SET}(B);
 \mathsf{NIP};\mathsf{DROP};\mathsf{NIP};\mathsf{MP};\mathsf{NIP}.
\]
After the premise blocks and SET, unless there has already been a halt,
the state is
\[
 (I(A,B),[A]S,B,0).
\]
The premise $A$ has remained below every temporary AX output.  The next
three commands give
\[
 (I(A,B),S,B,A)
 \longmapsto (A,[I(A,B)]S,B,A)
 \longmapsto (A,S,B,I(A,B)).
\]
If $A=1$, the last of these NIPs halts instead, which is an allowed outcome.
Otherwise the parameter-free MP succeeds and produces $(B,[0]S,B,0)$.
The final NIP removes that success zero, giving $(B,S,B,0)$.  No read of an
MP parameter is needed.  This argument uses the original success behavior
that pushes zero; an omit-zero MP would omit the final cleanup NIP as well.

This completes a finite structural induction on native derivations.  It
does not presume any substitution rule, well-formedness test, or decoding
of arbitrary numerical formula codes.

\subsection{Soundness, enumeration, and the refutation equivalence}

The soundness invariant is the accepted one: the current code, the cached
code, and every theorem-stack entry are native derivable.  The numeric
parameter $d$ has no theoremhood requirement.  DROP and NIP preserve this
invariant.  MP either returns truth or applies native modus ponens to its
two invariant premises.  AX emits native axioms and a generalization of
the old cached theorem; its gated output is either truth or an admissible
native bundle.  Its arbitrary residual changes only the unconstrained
parameter.  Consequently every halt, which occurs only upon NIP observing
cached code $1$, yields a native refutation.

Conversely, compile a native proof of $1$.  By the universal compiler
contract, either its block already halts or it returns $(1,S,d',0)$.
Appending $\mathsf{DROP};\mathsf{NIP};\mathsf{NIP}$ then halts: the first
NIP places the current falsity code in the cache, and the second detects it.
Thus there is a finite command word which halts from every work state.

The accepted exhaustive candidate shell still applies with the shorter
command grammar.  For a continuation code $z$, pack one command as
\[
 \begin{array}{c|c}
 \mathsf{MP}&P(0,z)\\
 \mathsf{DROP}&P(1,z)\\
 \mathsf{NIP}&P(2,z)\\
 \mathsf{AX}(c,y,x)&P(c+3,P(y,P(x,z))).
 \end{array}
\]
More formally, define the list cursor by $E([])=0$ and
$E(c::p)=1+\operatorname{pack}(c,E(p))$.  The natural number enumerated
for the nonempty list $c::p$ is $\operatorname{pack}(c,E(p))=E(c::p)-1$;
fetch loads this raw number at candidate entry, and decrements the positive
cursor before subsequent commands.  Thus a final MP with raw packed value
$P(0,0)=0$ is represented by cursor $1$, not by an empty cursor.  Each
command in an in-progress candidate strictly reduces
the positive cursor: fetch decrements it once, and subsequent Cantor tail
projections do not increase it.  Hence each candidate terminates or halts,
and the outer shell eventually reaches every later candidate number.
If bootstrap starts after some finite candidate index, prepend an arbitrary
AX with a sufficiently large matrix input to the universally halting word;
it remains universally halting, and its encoded candidate can be made
arbitrarily large.  This establishes completeness independently of which
finite initial candidates bootstrap skips.

Therefore the exhaustive command interpreter, started from a theorem-safe
seed with finite candidate index, halts if and only if the unchanged
native numerical calculus derives $1$.  A theorem-safe seed means that
its current value, cache, and stack entries are native derivable; it places
no restriction on the numeric parameter.  The zero work state, or the
accepted bootstrap's theorem-safe work state, meets this condition.
Arbitrary initial work states are used only in the universal completeness
contract, not in the soundness assertion.  For example, an arbitrary cache
already containing $1$ would trivially halt on NIP.
Composing with the repository's
accepted native-calculus calibration gives the same ZFC consistency
comparison.  All constructions above are finite syntactic transformations
and elementary integer identities, so their mathematical proofs are
formalizable in ZFC.  A new exact-table mechanization has not been supplied
by this command-level argument.

\subsection{Source anchors and verification scope}

The raw stack alias $[0]0=P(0,0)=0$ causes no exception in these proofs:
every stack step uses the exact identity $P^{-1}(P(x,S))=(x,S)$, which
also holds for $x=S=0$.  Displayed list lengths count the planned finite
pops, without assuming that the numerical stack has a separate end marker.

The older Lean comparison computes an absolute distance internally; that
internal description must not be substituted for the one-sided residual
of the optimized literal K217 gate.  Both comparisons implement the same
destructive equality interface before this modification.

The inherited register and macro contracts are in
\nolinkurl{L2Lib254.lean}, \nolinkurl{L2Macros254.lean},
\nolinkurl{L2Control254.lean}, \nolinkurl{L2Axioms254.lean}, and
\nolinkurl{L2Root254.lean}, under
\nolinkurl{Computability/BusyBeaver/Lean/BusyBeaver/ZFC/} at repository commit
\nolinkurl{d3d6b8c14e1ddff37eb1417c5448caa41cfd6da7}.
In particular, \nolinkurl{loop124}, \nolinkurl{ifeq_ne}, and
\nolinkurl{ifeq_clear} identify the exact destructive comparison;
\nolinkurl{l71} composes the two gates.  The stack order and native compiler
comparison are given by \nolinkurl{Enum222.lean}'s \nolinkurl{axList},
\nolinkurl{axStack_eq}, \nolinkurl{select_entry}, and \nolinkurl{build}.
The proof above is a new mathematical adaptation of those interfaces, not
a claim that those existing Lean files already verify the revised program.

\paragraph{Relation to later upstream residual-counter research.}
While this investigation was in progress, upstream intake 193 published
Source153, including \nolinkurl{theory/counter_salt.tex}.  That work already
establishes the useful principle that the parameter register need not hold
a theorem and may retain an arithmetic comparison residual.  Its repaired
reader transfers the residual before command decoding; AX leaves
$\max(c-a-1,0)$, and failed MP leaves
$a+\max(I(t,a)-h-1,0)$, where $a=d+x$.  It retains an explicit MP input
$x$, and its compiler restores the former zero-parameter interface.

The present contribution should therefore not be described as the invention
of residual-counter freedom.  The distinct step is a \emph{parameter-free
MP}: an AX setter supplies any desired consequent while preserving the
current implication and protecting the other premise below the temporary
AX stack prefix.  The theorem above uses an absolute AX input, consumes
the old $U$ only after its final use, and permits a nonzero parameter at
compiler-block exits.  These choices differ from Source153's reported
contract.  Neither the existing residual freedom nor this new logical
compilation result alone establishes a smaller binary transition table.

For exact attribution, Source153 arrived in ProveIt commit
\begin{center}\footnotesize\ttfamily
1e87618cf2c291cae16a66724f5c76b271e90c45
\end{center}
as \nolinkurl{docs/incoming/ZFC_218_Research_2026-10-20 (1).zip}.
The retained reference is the archive member
\nolinkurl{ZFC_218_Research_2026-10-20/theory/counter_salt.tex},
whose exact SHA-256 is
\begin{center}\footnotesize\ttfamily
19571441839a01021193dd25db163576ada1752b8ba20acf0b858a061f1b127c.
\end{center}
The 11,715-byte original TeX, its attribution and license, and the precise
archive/member hashes and URL are preserved in the supplied semantic
references, with provenance in \nolinkurl{SOURCE153_PINS.json}.

The accompanying \nolinkurl{check_residual_contracts.py} passes 4,619 exact
integer cases for the helper contracts, all row-selection positions,
the generalization and MP induction steps, terminal falsity detection, and
command packing, including 432 additional opposite-residual setter cases.
It uses three permitted instantiations of the AX
safe-row interface, including all-zero rows, and tests protected stacks
containing falsity as well as the zero-stack alias.  Premise-block test
oracles implement the induction hypothesis and are explicitly not machine
instructions.  Consequently this checker is a focused check of stack and
command algebra; it does not evaluate the complete native emitter, prove
the universal theorem, or establish the binary compilation.
